\section{Данные}

\begin{frame}{Данные: сколько и откуда}
  Объём данных не заменяет качество и разнообразие примеров.
  \begin{description}
    \item[Открытые датасеты] Готовые изображения, тексты и таблицы.
      Проверьте лицензию и соответствие задаче.
    \item[Краудсорсинг] Люди размечают данные по инструкции.
      Часть примеров проверяют несколько раз.
    \item[Покупка] Для некоторых задач набор можно купить; для других его придётся собрать.
    \item[Без разметки] Учебную задачу можно построить из самих данных:
      например, восстанавливать скрытую часть текста.
  \end{description}
\end{frame}

\begin{frame}{Типы признаков}
  Строки таблицы --- объекты, столбцы --- признаки.
  \medskip\par
  \begin{tabularx}{\linewidth}{@{}l Y@{}}
    \toprule
    \textbf{Тип} & \textbf{Значения и примеры}\\
    \midrule
    Численный & Число: рост, доход. Бывает вещественным или целочисленным.\\
    Бинарный & Два значения: да или нет, 0 или 1. Можно рассматривать как число или категорию.\\
    Категориальный & Вариант из набора: профессия, день недели.\\
    Ординальный & Категория с порядком: класс опасности от 1 до 4, год обучения.\\
    \bottomrule
  \end{tabularx}
\end{frame}

\begin{frame}[c]{Какого типа признак?}
  \small Численный $\cdot$ Категориальный $\cdot$ Бинарный $\cdot$ Ординальный

  \smallskip
  \begin{columns}[T]
    \begin{column}{0.47\textwidth}
      \begin{numrows}
        1 & Рост человека\\
        2 & Профессия\\
        3 & Оплатил ли клиент заказ\\
      \end{numrows}\par
    \end{column}
    \begin{column}{0.47\textwidth}
      \begin{numrows}
        4 & Класс энергопотребления (от A до G)\\
        5 & День недели\\
        6 & Год обучения студента\\
      \end{numrows}\par
    \end{column}
  \end{columns}
\end{frame}

\begin{frame}[c]{Ответ: типы признаков}
  \small
  \zebra{2}
  \begin{tabular}{>{\raggedright\arraybackslash}p{0.444\textwidth}l}
    \toprule
    \textbf{Признак} & \textbf{Тип}\\
    \midrule
    Рост человека & \textbf{Численный}\\
    Профессия & \textbf{Категориальный}\\
    Оплатил ли заказ & \textbf{Бинарный}\\
    Класс энергопотребления, A--G & \textbf{Ординальный}\\
    День недели & \textbf{Категориальный}\\
    Год обучения & \textbf{Ординальный}\\
    \bottomrule
  \end{tabular}
\end{frame}

\begin{frame}[c]{Что не так с этими данными?}

  \medskip
  \large
  \zebra{2}\renewcommand{\arraystretch}{1.7}\setlength{\tabcolsep}{12pt}
  \begin{tabular}{llrl}
    \toprule
    \textbf{Клиент} & \textbf{Возраст} & \textbf{Доход, руб.} & \textbf{Вернул кредит}\\
    \midrule
    Анна    & 34  & 90\,000    & да\\
    Борис   & N/A & 70\,000    & нет\\
    Виктор  & 140 & 85\,000    & да\\
    Галина  & 29  & 2\,400\,000 & да\\
    Дмитрий & 41  & 65\,000    & да\\
    \bottomrule
  \end{tabular}

  \bigskip
  \textbf{Часть проблем в таблице не видна вовсе. Какие это?}
\end{frame}

\begin{frame}{Ответ: четыре типа проблем в данных}
  \begin{description}
    \item[Пропуски] Возраст Бориса неизвестен. Запись можно удалить,
      значение заполнить или сохранить пропуск как признак.
    \item[Выбросы] Возраст 140 лет требует проверки.
      Большой доход может быть реальным; необычные значения не удаляют автоматически.
    \item[Ошибки разметки] В таблице могут быть неверные ответы,
      например из-за неоднозначной инструкции разметчикам.
    \item[Изменение данных] Новые данные могут отличаться от обучающих.
      Изменение распределения называют data drift.
  \end{description}
\end{frame}
